% DPLL report
% Tables under tables/ are generated by tests/make_tables_dpll.py from
% results/task2*.csv, so `make bench` refreshes them

\documentclass[11pt,a4paper]{article}
% Shared preamble for the NAIL094 reports
% Used by every task via % Shared preamble for the NAIL094 reports
% Used by every task via % Shared preamble for the NAIL094 reports
% Used by every task via \input{../../shared/preamble}

\usepackage[T1]{fontenc}
\usepackage[utf8]{inputenc}
\usepackage{lmodern}
\usepackage[margin=2.5cm]{geometry}

\usepackage{amsmath}
\usepackage{amssymb}
\usepackage{booktabs}
\usepackage{float}
\usepackage{graphicx}
\usepackage{listings}
\usepackage{xcolor}
\usepackage{caption}

\usepackage[hidelinks]{hyperref}

\setlength{\parskip}{0.5em}
\setlength{\parindent}{0pt}

% Title fields, set these in the report before \begin{document}
\makeatletter
\newcommand{\reportcourse}[1]{\def\@reportcourse{#1}}
\newcommand{\reporttask}[1]{\def\@reporttask{#1}}
\newcommand{\reportauthor}[1]{\def\@reportauthor{#1}}

\newcommand{\makereporttitle}{%
  \begin{center}
    {\LARGE\bfseries \@reportcourse \par}
    \vspace{0.8em}
    {\Large \@reporttask \par}
    \vspace{0.6em}
    {\normalsize \@reportauthor \par}
  \end{center}
  \vspace{1em}
  \hrule
  \vspace{1.5em}
}
\makeatother

% Inline code
\newcommand{\code}[1]{\texttt{#1}}

% Code listings
\lstdefinestyle{reportcode}{
  basicstyle=\ttfamily\small,
  keywordstyle=\color{blue!60!black},
  commentstyle=\color{gray},
  stringstyle=\color{green!40!black},
  breaklines=true,
  showstringspaces=false,
  frame=single,
  rulecolor=\color{gray!40},
  numbers=none,
}
\lstset{style=reportcode}


\usepackage[T1]{fontenc}
\usepackage[utf8]{inputenc}
\usepackage{lmodern}
\usepackage[margin=2.5cm]{geometry}

\usepackage{amsmath}
\usepackage{amssymb}
\usepackage{booktabs}
\usepackage{float}
\usepackage{graphicx}
\usepackage{listings}
\usepackage{xcolor}
\usepackage{caption}

\usepackage[hidelinks]{hyperref}

\setlength{\parskip}{0.5em}
\setlength{\parindent}{0pt}

% Title fields, set these in the report before \begin{document}
\makeatletter
\newcommand{\reportcourse}[1]{\def\@reportcourse{#1}}
\newcommand{\reporttask}[1]{\def\@reporttask{#1}}
\newcommand{\reportauthor}[1]{\def\@reportauthor{#1}}

\newcommand{\makereporttitle}{%
  \begin{center}
    {\LARGE\bfseries \@reportcourse \par}
    \vspace{0.8em}
    {\Large \@reporttask \par}
    \vspace{0.6em}
    {\normalsize \@reportauthor \par}
  \end{center}
  \vspace{1em}
  \hrule
  \vspace{1.5em}
}
\makeatother

% Inline code
\newcommand{\code}[1]{\texttt{#1}}

% Code listings
\lstdefinestyle{reportcode}{
  basicstyle=\ttfamily\small,
  keywordstyle=\color{blue!60!black},
  commentstyle=\color{gray},
  stringstyle=\color{green!40!black},
  breaklines=true,
  showstringspaces=false,
  frame=single,
  rulecolor=\color{gray!40},
  numbers=none,
}
\lstset{style=reportcode}


\usepackage[T1]{fontenc}
\usepackage[utf8]{inputenc}
\usepackage{lmodern}
\usepackage[margin=2.5cm]{geometry}

\usepackage{amsmath}
\usepackage{amssymb}
\usepackage{booktabs}
\usepackage{float}
\usepackage{graphicx}
\usepackage{listings}
\usepackage{xcolor}
\usepackage{caption}

\usepackage[hidelinks]{hyperref}

\setlength{\parskip}{0.5em}
\setlength{\parindent}{0pt}

% Title fields, set these in the report before \begin{document}
\makeatletter
\newcommand{\reportcourse}[1]{\def\@reportcourse{#1}}
\newcommand{\reporttask}[1]{\def\@reporttask{#1}}
\newcommand{\reportauthor}[1]{\def\@reportauthor{#1}}

\newcommand{\makereporttitle}{%
  \begin{center}
    {\LARGE\bfseries \@reportcourse \par}
    \vspace{0.8em}
    {\Large \@reporttask \par}
    \vspace{0.6em}
    {\normalsize \@reportauthor \par}
  \end{center}
  \vspace{1em}
  \hrule
  \vspace{1.5em}
}
\makeatother

% Inline code
\newcommand{\code}[1]{\texttt{#1}}

% Code listings
\lstdefinestyle{reportcode}{
  basicstyle=\ttfamily\small,
  keywordstyle=\color{blue!60!black},
  commentstyle=\color{gray},
  stringstyle=\color{green!40!black},
  breaklines=true,
  showstringspaces=false,
  frame=single,
  rulecolor=\color{gray!40},
  numbers=none,
}
\lstset{style=reportcode}


\reportcourse{NAIL094 --- Decision Procedures and SAT/SMT Solvers}
\reporttask{DPLL Experimental Evaluation}
\reportauthor{Zerina Ahmetovic}

\begin{document}
\makereporttitle

\section{Task overview}

The implemented DPLL solver with basic branch and bound DPLL algorithm and unit propagation is
tested on small benchmarks. The correct input format is either SMT-LIB or DIMACS, and output writes
out whether the formula is SAT (in which case a model is shown as well) or UNSAT, along with statistical
information (CPU time, number of decisions and steps of unit propagation).

Two sets of benchmarks are used. The first is the five example formulas given
with the preceding task of the course~\cite{task1}, which asks for a Tseitin
encoding of NNF formulas written in a simplified SMT-LIB syntax
(Table~\ref{tab:examples}). The second is three families from
SATLIB~\cite{satlib}: All Interval Series~\cite{ais}, pigeonhole~\cite{phole}
and Uniform Random-3-SAT~\cite{uf}, collected in Table~\ref{tab:satlib}.
The first two are structured and are the
families suggested for this task; \texttt{ais6}, \texttt{ais8}, \texttt{ais10},
\texttt{hole6} and \texttt{hole7} are all included, together with the larger
\texttt{ais12}, \texttt{hole8} and \texttt{hole9} so that the growth in solving
time is visible over a wider range. The random instances are included for
contrast only, being uniformly generated rather than built from a combinatorial
principle.

\section{Implementation}

The solver is invoked as \texttt{dpll [input]}. The input format follows the
file extension, \texttt{.cnf} for DIMACS and \texttt{.sat} for the simplified
SMT-LIB format, and can be forced with \texttt{--cnf} or \texttt{--sat}.
SMT-LIB input is passed through the Tseitin encoder written for the preceding
task~\cite{task1}, so the solver itself only ever sees CNF.

The algorithm consists of three fundamental parts, explained below.

\textbf{Trail: } The current partial assignment, kept three ways at once: a value
per literal so that testing whether a literal is false is a single index flip, a
stack in assignment order so that undoing runs in reverse, and a mark per
decision level recording where that level began. The decision literal of level
$d+1$ is the first entry after mark $d$, which is how a decision is recovered
when it has to be flipped. The stack doubles as the propagation queue.

\textbf{Unit propagation:} For each literal $l$, \texttt{occ[l]} lists the
clauses containing it. Assigning $l$ true can only shorten clauses containing
$\lnot l$, so only those are visited; the work per assignment is proportional to
the occurrences of the assigned literal rather than to the size of the formula.

\textbf{Search:} Branch on the lowest-indexed unassigned variable, positive
first, which the task permits. On a conflict the solver walks up to the deepest
decision level that has not yet been tried in both polarities, discards
everything below it and re-decides the negated literal; if no such level
remains the formula is unsatisfiable. One bit per decision level replaces the
recursion, so the search is a flat loop.

Clauses are normalised once at load: repeated literals collapsed, tautological
clauses dropped, and clauses repeating an earlier one dropped. The last matters
on All Interval Series, where \texttt{ais6} loses 140 of its 581 clauses as
duplicates, which is why Table~\ref{tab:satlib} reports 441 clauses for it
rather than the 581 of the original file.

\section{Experimental evaluation}

The times in Tables~\ref{tab:examples} and~\ref{tab:satlib} were measured with
g++ 14.2.0 (MSYS2 UCRT64) at \texttt{-std=c++20 -O2}. Each instance was run
three times and the fastest kept, the timer covers the search only, reading the
file or encoding it are not included.

\subsection*{Example formulas in SMT-LIB format}

\begin{table}[H]
\centering
\small
\begin{tabular}{lrrlrrr}
\toprule
instance & vars & clauses & result & decisions & unit props & CPU s \\
\midrule
toy\_5.sat & 9 & 13 & SAT & 5 & 3 & 0.000 \\
toy\_50.sat & 99 & 148 & SAT & 50 & 48 & 0.000 \\
toy\_100.sat & 199 & 298 & SAT & 100 & 98 & 0.000 \\
nested\_5.sat & 22 & 45 & SAT & 6 & 15 & 0.000 \\
nested\_8.sat & 136 & 371 & SAT & 65 & 1\,596 & 0.000 \\
\bottomrule
\end{tabular}
\caption{The five example inputs from the first task, read in the simplified SMT-LIB format and encoded by Tseitin before solving. Variable and clause counts are those of the resulting CNF.}
\label{tab:examples}
\end{table}


As Table~\ref{tab:examples} shows, all five examples are satisfiable, and only
\texttt{nested\_8} produces any conflicts at all (31). The other four are solved
without backtracking once, so the Tseitin
encoding of a satisfiable NNF formula leaves the solver very little to do.

\subsection*{SATLIB benchmarks}

\begin{table}[H]
\centering
\small
\begin{tabular}{lrrlrrr}
\toprule
instance & vars & clauses & result & decisions & unit props & CPU s \\
\midrule
hole6 & 42 & 133 & UNSAT & 6\,490 & 29\,322 & 0.001 \\
hole7 & 56 & 204 & UNSAT & 65\,560 & 318\,495 & 0.009 \\
hole8 & 72 & 297 & UNSAT & 756\,686 & 3\,933\,538 & 0.124 \\
hole9 & 90 & 415 & UNSAT & 9\,825\,028 & 54\,417\,935 & 1.635 \\
ais6 & 61 & 441 & SAT & 57 & 515 & 0.000 \\
ais8 & 113 & 1149 & SAT & 550 & 7\,423 & 0.001 \\
ais10 & 181 & 2377 & SAT & 8\,973 & 138\,777 & 0.017 \\
ais12 & 265 & 4269 & SAT & 192\,780 & 3\,389\,633 & 0.355 \\
uf50-01 & 50 & 218 & SAT & 1\,016 & 7\,573 & 0.001 \\
uf50-010 & 50 & 218 & SAT & 370 & 2\,526 & 0.001 \\
uf50-0100 & 50 & 218 & SAT & 318 & 2\,355 & 0.000 \\
uf50-0101 & 50 & 218 & SAT & 44 & 224 & 0.000 \\
\bottomrule
\end{tabular}
\caption{SATLIB instances from three categories. CPU time is the fastest of three runs and covers search only, not parsing. Clause counts are after duplicate and tautological clauses are removed at load.}
\label{tab:satlib}
\end{table}


Table~\ref{tab:satlib} shows that inside the pigeonhole family the decision
count grows by factors of $10.1$,
$11.5$ and $13.0$ from \texttt{hole6} to \texttt{hole9}, so the growth factor
itself increases rather than staying constant. CPU time follows it, from
$0.001$\,s to $1.635$\,s for 48 additional variables. A running time of
$2^{cn}$ would give a constant factor; one rising by about $1.4$ per step is
the shape of $n!$-type growth, that is $2^{\Theta(n \log n)}$. That is the
expected shape here, since DPLL without clause learning corresponds to
tree-like resolution, for which the pigeonhole principle is known to require
refutations of that size.

Size alone does not predict difficulty. Across families in the same table,
\texttt{ais12} has 265
variables against \texttt{hole9}'s 90 and is still solved in $0.355$\,s rather
than $1.635$\,s, being satisfiable and far less constrained. Within a family the
same holds: the four \texttt{uf50} instances are identical in size yet differ by
a factor of $23$ in decisions. These sit at the phase-transition ratio
$218/50 = 4.36$, where difficulty is heavy-tailed even at fixed size, so with a
fixed index-order decision rule how soon the search meets a solution varies
enormously from one instance to the next. Size alone predicts neither.

The task also asks for the number of unit propagation steps, which
Table~\ref{tab:satlib} reports. Read per decision it separates the families:
pigeonhole stays near $4.5$ to $5.5$ and \texttt{uf50} near $5$ to $7.5$, while
All Interval Series climbs from $9.0$ on \texttt{ais6} to $17.6$ on
\texttt{ais12}. Each decision there forces far more assignments, and
increasingly so with size, which is the implication structure the later
look-ahead solver exploits.

\subsection*{Code}

\begin{sloppypar}
The full source is in the accompanying repository~\cite{repo}.
\texttt{src/dpll.hpp} holds the trail, the statistics, the adjacency-list
propagation and the search loop; \texttt{src/dpll.cpp} is the command line,
format dispatch and output. The DIMACS reader is \texttt{src/cnf.hpp} and the
SMT-LIB path reuses \texttt{src/formula.hpp} and \texttt{src/tseitin.hpp},
written for the encoding task~\cite{task1}. Correctness is checked by
\texttt{tests/check\_solver.py},
which compares the verdict against exhaustive search and validates the reported
model on both input formats, and by \texttt{tests/fuzz.py} over random
instances. Measurements come from \texttt{tests/bench\_dpll.py}.
\end{sloppypar}

\subsection*{Summary}

The solver reads both required formats, reports SAT or UNSAT with a model, and
prints the three required statistics. Solving time grows steeply with instance
size within a family, most sharply on the unsatisfiable pigeonhole instances,
but the comparison across families in Table~\ref{tab:satlib} shows that
structure and satisfiability matter more than size.

\begin{thebibliography}{9}

\bibitem{repo}
Z.~Ahmetovic.
\newblock Source repository for the NAIL094 programming tasks.
\newblock \url{https://github.com/inferno73/SAT-solver-NAIL094}

\bibitem{task1}
P.~Ku\v{c}era.
\newblock Tseitin encoding and DIMACS format, programming task for NAIL094.
\newblock \url{https://ktiml.mff.cuni.cz/~kucerap/satsmt/practical/task_tseitin.php}

\bibitem{satlib}
H.~H.~Hoos and T.~St\"utzle.
\newblock SATLIB: An Online Resource for Research on SAT.
\newblock \url{https://www.cs.ubc.ca/~hoos/SATLIB/benchm.html}

\bibitem{ais}
SATLIB benchmark problems, All Interval Series.
\newblock \url{https://www.cs.ubc.ca/~hoos/SATLIB/Benchmarks/SAT/AIS/descr.html}

\bibitem{phole}
SATLIB benchmark problems, DIMACS pigeon hole.
\newblock \url{https://www.cs.ubc.ca/~hoos/SATLIB/Benchmarks/SAT/DIMACS/PHOLE/descr.html}

\bibitem{uf}
SATLIB benchmark problems, Uniform Random-3-SAT.
\newblock \url{https://www.cs.ubc.ca/~hoos/SATLIB/Benchmarks/SAT/RND3SAT/descr.html}

\end{thebibliography}

\end{document}
